%=====================================================================
\chapter{Inferential Statistics}\label{ch:inference}
%=====================================================================
\locator{2}{Part II --- From Data to Insight}

\begin{outcomes}
By the end of this chapter you will be able to:
\begin{tight}
  \item \textbf{Distinguish} a population from a sample and identify the sampling method used in a study.
  \item \textbf{Construct} and \textbf{interpret} a confidence interval, and state precisely what it does \emph{not} mean.
  \item \textbf{Conduct} a hypothesis test and interpret a $p$-value correctly.
  \item \textbf{Explain} why statistical significance and practical importance are different things, using effect size.
  \item \textbf{Recognise} the multiple-comparisons problem and apply a correction.
\end{tight}
\end{outcomes}

\begin{prereq}
\chref{probability} --- especially the Central Limit Theorem, which is what licenses
everything in this chapter --- and \chref{descriptive}.
\end{prereq}

\begin{keyterms}
population $\bullet$ sample $\bullet$ sampling bias $\bullet$ standard error
$\bullet$ confidence interval $\bullet$ null and alternative hypothesis $\bullet$
$p$-value $\bullet$ significance level $\bullet$ Type I and Type II error $\bullet$
statistical power $\bullet$ effect size $\bullet$ multiple comparisons $\bullet$
$p$-hacking
\end{keyterms}

\section{From Sample to Population}

\begin{definitionbox}[Population and Sample]
The \term{population} is every unit you want to make a claim about. The
\term{sample} is the subset you actually measured. Inference is the discipline of
saying what the sample entitles you to claim about the population --- and, equally,
what it does not.
\end{definitionbox}

\rowcolors{2}{white}{lightbg}
\begin{longtable}{@{}p{3.0cm}p{5.2cm}p{6.3cm}@{}}
\toprule
\rowcolor{primary!15}
\textbf{Sampling method} & \textbf{How} & \textbf{Bias risk} \\
\midrule
\endhead
Simple random & Every unit equally likely & Low, but often impractical \\
Stratified & Sample within each subgroup & Low; \emph{improves} precision when groups differ \\
Cluster & Sample whole groups (e.g.\ 5 classes) & Moderate; units within a cluster are similar, so effective $n$ is smaller than it looks \\
Systematic & Every $k$-th unit & Low unless the data has a matching period \\
Convenience & Whoever is available & \textbf{High} --- the commonest and most damaging method in student projects \\
Voluntary response & Whoever chooses to reply & \textbf{High} --- respondents differ systematically from non-respondents \\
\bottomrule
\end{longtable}

\begin{alertbox}[No Sample Size Fixes a Biased Sample]
If your survey reaches only students who log in to the VLE, no increase in $n$
will tell you anything about students who have stopped logging in --- and in a
dropout study, those are the population of interest. A large biased sample is
\emph{more} dangerous than a small one, because its narrow confidence interval
projects false authority.
\end{alertbox}

\section{Confidence Intervals}

\begin{definitionbox}[Standard Error and Confidence Interval]
The \term{standard error} of the mean is $\text{SE} = s/\sqrt{n}$ --- the typical
distance between a sample mean and the true mean. An approximate 95\%
\term{confidence interval} is
\[
\bar{x} \;\pm\; 1.96 \times \frac{s}{\sqrt{n}}
\]
\end{definitionbox}

\begin{worked}[A Confidence Interval, and What It Means]
A sample of $n = 100$ support tickets has mean resolution time
$\bar{x} = 42$ minutes with $s = 15$ minutes.

\smallskip
\textbf{Standard error:} $\text{SE} = 15/\sqrt{100} = 15/10 = 1.5$ minutes.

\textbf{Margin of error:} $1.96 \times 1.5 = 2.94$ minutes.

\textbf{95\% CI:} $42 \pm 2.94 = [39.1,\; 44.9]$ minutes.

\smallskip
\textbf{What this means:} if you repeated this sampling procedure many times, about
95\% of the intervals so constructed would contain the true population mean.

\textbf{What it does \emph{not} mean:} that there is a 95\% probability the true
mean lies in \emph{this} interval. The true mean is a fixed number; it either is or
is not in $[39.1, 44.9]$. The 95\% describes the reliability of the
\emph{procedure}, not this one result.

\smallskip
\textbf{Now quadruple the sample.} With $n = 400$, $\text{SE} = 15/20 = 0.75$ and
the interval halves to $[40.5, 43.5]$. Four times the data for twice the precision
--- the $\sqrt{n}$ from \chref{probability}, in its practical form. Budget
accordingly.
\end{worked}

\dsfig{%
\begin{tikzpicture}
\begin{axis}[
  width=12.2cm, height=5.4cm,
  xlabel={\scriptsize sample size $n$}, ylabel={\scriptsize margin of error (min)},
  xmin=0, xmax=820, ymin=0, ymax=11,
  axis lines=left, grid=major, grid style={gray!18},
  tick label style={font=\tiny}, label style={font=\scriptsize},
]
\addplot[primary, line width=1.4pt, domain=10:800, samples=200] {1.96*15/sqrt(x)};
\addplot[accent, only marks, mark=*, mark size=2pt] coordinates {(100,2.94) (400,1.47)};
\draw[gray!60, dashed] (axis cs:100,0) -- (axis cs:100,2.94) -- (axis cs:0,2.94);
\draw[gray!60, dashed] (axis cs:400,0) -- (axis cs:400,1.47) -- (axis cs:0,1.47);
\node[font=\tiny, accent, anchor=south west] at (axis cs:105,3.1) {$n=100$: $\pm2.94$};
\node[font=\tiny, accent, anchor=south west] at (axis cs:405,1.6) {$n=400$: $\pm1.47$};
\node[font=\scriptsize\itshape, text=gray!55!black, align=center, anchor=north east]
     at (axis cs:790,9.5) {diminishing returns:\\to halve the error,\\quadruple the data};
\end{axis}
\end{tikzpicture}}%
{Margin of error against sample size. The curve is $1.96s/\sqrt{n}$. The flat right-hand
region is why collecting ten times more data rarely improves a study ten-fold --- and why
reducing bias usually beats increasing $n$.}{margin-error}

\section{Hypothesis Testing}

\begin{definitionbox}[Null and Alternative Hypotheses]
The \term{null hypothesis} $H_0$ states there is no effect --- no difference, no
relationship. The \term{alternative} $H_1$ states there is one. A test asks: if
$H_0$ were true, how surprising would data like mine be?
\end{definitionbox}

\begin{definitionbox}[$p$-value]
The $p$-value is the probability of observing data \emph{at least as extreme as
yours}, assuming the null hypothesis is true. Small $p$ means your data would be
unusual if there were no effect.
\end{definitionbox}

\begin{alertbox}[What a $p$-value Is Not]
It is \textbf{not} the probability that $H_0$ is true.\\
It is \textbf{not} the probability your result was due to chance.\\
It is \textbf{not} a measure of how large or important the effect is.\\
$p = 0.04$ and $p = 0.0001$ do not mean ``true'' and ``very true''.
\end{alertbox}

\dsfig{%
\begin{tikzpicture}[font=\scriptsize]
\node[rounded corners=3pt, fill=primary, text=white, font=\scriptsize\bfseries,
      minimum width=3.4cm, minimum height=0.7cm] at (4.2,1.55) {THE TRUTH (unknown)};
\node[font=\scriptsize\bfseries, text=primary] at (1.0,0.85) {$H_0$ is true};
\node[font=\scriptsize\bfseries, text=primary] at (7.4,0.85) {$H_0$ is false};
\node[rounded corners=3pt, fill=primary!75, text=white, font=\scriptsize\bfseries,
      minimum width=1.0cm, minimum height=3.0cm, rotate=90] at (-3.0,-1.35) {YOUR DECISION};

\node[align=center, font=\scriptsize\bfseries, text=primary, text width=2.2cm] at (-1.4,-0.3)
     {do not reject $H_0$};
\node[align=center, font=\scriptsize\bfseries, text=primary, text width=2.2cm] at (-1.4,-2.4)
     {reject $H_0$};

\node[rounded corners=3pt, fill=greenhl!14, draw=greenhl, text width=5.2cm, align=center,
      minimum height=1.7cm, inner sep=4pt] at (1.0,-0.3)
     {{\color{greenhl!50!black}\bfseries Correct}\\[2pt]
      \emph{no effect, none claimed}};
\node[rounded corners=3pt, fill=accent!12, draw=accent, text width=5.2cm, align=center,
      minimum height=1.7cm, inner sep=4pt] at (7.4,-0.3)
     {{\color{accent}\bfseries TYPE II ERROR ($\beta$)}\\[2pt]
      \emph{a real effect missed}\\
      {\color{accent!70!black}\tiny e.g.\ an intrusion not detected}};
\node[rounded corners=3pt, fill=accent!12, draw=accent, text width=5.2cm, align=center,
      minimum height=1.7cm, inner sep=4pt] at (1.0,-2.4)
     {{\color{accent}\bfseries TYPE I ERROR ($\alpha$)}\\[2pt]
      \emph{an effect claimed that is not there}\\
      {\color{accent!70!black}\tiny e.g.\ a false alarm}};
\node[rounded corners=3pt, fill=greenhl!14, draw=greenhl, text width=5.2cm, align=center,
      minimum height=1.7cm, inner sep=4pt] at (7.4,-2.4)
     {{\color{greenhl!50!black}\bfseries Correct --- POWER ($1-\beta$)}\\[2pt]
      \emph{a real effect detected}};

\node[anchor=north, align=left, text width=13.5cm, font=\scriptsize, text=gray!55!black] at (3.2,-4.0)
     {Setting $\alpha = 0.05$ means accepting a 1-in-20 Type I error rate \emph{by design}.
      Lowering $\alpha$ reduces false alarms but raises misses --- the same trade-off that
      reappears as precision versus recall in \chref{evaluation}.};
\end{tikzpicture}}%
{The two errors. Which one is worse is a domain question, never a statistical one: a Type I
error wastes an analyst's afternoon, a Type II error lets an intrusion through.}{error-types}

\begin{worked}[Did the New Lab Format Help?]
Pass rates: old format 68 of 100 students; new format 79 of 100.

\smallskip
\textbf{$H_0$:} the true pass rates are equal. \quad \textbf{$H_1$:} they differ.

\textbf{Test:} two-proportion z-test.
Pooled proportion $\hat{p} = (68+79)/200 = 0.735$.
\[
\text{SE} = \sqrt{0.735 \times 0.265 \times \left(\tfrac{1}{100}+\tfrac{1}{100}\right)}
          = \sqrt{0.1948 \times 0.02} = 0.0624
\]
\[
z = \frac{0.79 - 0.68}{0.0624} = \frac{0.11}{0.0624} = 1.76
\]

\textbf{$p$-value:} for a two-sided test, $p \approx 0.078$.

\smallskip
\textbf{Decision at $\alpha = 0.05$:} do not reject $H_0$. The evidence is
suggestive but not conclusive.

\textbf{The honest report:} ``Pass rates rose by 11 percentage points
(68\% to 79\%). This did not reach significance at the 5\% level ($p = 0.078$),
which with 100 students per group is unsurprising --- the study had roughly 40\%
power to detect an effect of this size. The observed difference is educationally
meaningful and warrants a larger trial, not dismissal.''

\smallskip
\textbf{The dishonest report:} ``No significant difference; the new format does not
work.'' Absence of evidence is not evidence of absence, and with this sample size
the test was never capable of settling the question.
\end{worked}

\section{Effect Size and Power}

\begin{definitionbox}[Effect Size]
A measure of \emph{how big} a difference is, independent of sample size. For two
means, Cohen's $d = \dfrac{\bar{x}_1 - \bar{x}_2}{s_{\text{pooled}}}$, conventionally
read as small ($\approx 0.2$), medium ($\approx 0.5$), large ($\approx 0.8$).
\end{definitionbox}

\begin{conceptbox}[Why Every Result Needs Both Numbers]
With $n = 10$ million, a difference of 0.01\% in click rate will be highly
significant and completely worthless. With $n = 20$, a difference that would
transform the department may not reach significance at all. The $p$-value tells you
whether the effect is \emph{detectable}; the effect size tells you whether it is
\emph{worth acting on}. Report both, always.
\end{conceptbox}

\section{Multiple Comparisons}

\begin{alertbox}[The Problem]
At $\alpha = 0.05$, each test has a 1-in-20 chance of a false positive. Run 20
independent tests on data with no real effects and the probability of at least one
``significant'' result is $1 - 0.95^{20} = 64\%$. Run 100 and it is 99.4\%.

This is how a team ``discovers'' that students born in March do better: they tested
enough hypotheses that something had to come up.
\end{alertbox}

\begin{definitionbox}[Bonferroni Correction]
When running $m$ tests, use $\alpha' = \alpha / m$ for each. With 20 tests at an
overall 5\% level, require $p < 0.0025$ individually. It is conservative --- it
raises Type II errors --- but it is simple and defensible.
\end{definitionbox}

\begin{pitfall}[$p$-hacking]
\begin{tight}
  \item Testing many outcomes and reporting only the significant ones.
  \item Adding data until $p$ drops below 0.05, then stopping.
  \item Trying several subgroups and reporting the one that worked.
  \item Deciding the hypothesis after seeing the result.
\end{tight}
\textbf{The defence:} write the hypothesis, the test and the success criterion down
\emph{before} you look at the data --- exactly the discipline
\chref{lifecycle} demanded in framing, and which \chref{research} formalises as
pre-registration.
\end{pitfall}

%---------------------------------------------------------------------
\begin{fourlenses}{Inference in Practice}
\lensLA{Comparing pass rates between two teaching formats is a two-proportion
test --- but students are not independent: they share tutors, timetables and study
groups, which is cluster sampling in disguise, so the effective sample is smaller
than the headcount. Always report the effect size in percentage points alongside
$p$, because a statistically significant 0.5-point rise will not justify redesigning
a module.}

\lensPM{Sprint samples are small ($n \approx 10$--20), so power is low and
non-significant results are near-guaranteed. Confidence intervals on velocity are
far more useful to a team than tests: ``we complete 28--41 points per sprint with
95\% confidence'' supports planning in a way that ``$p = 0.31$'' never will.}

\lensIOT{With millions of readings, \emph{everything} is statistically significant
--- a $0.02^\circ$C difference between two sensor batches will have $p < 10^{-9}$
and mean nothing. In this domain the $p$-value is almost useless and the effect
size is almost everything. Beware also that consecutive readings are
autocorrelated, so the effective sample size is far below the row count.}

\lensSEC{Type I and Type II errors have wildly asymmetric costs here, so
$\alpha = 0.05$ is an arbitrary import from another field. Set the threshold from
the cost ratio instead (\chref{probability}'s expected-value calculation). Multiple
comparisons are acute: a SIEM testing thousands of rules per hour will generate
false positives continuously by construction unless corrected.}
\end{fourlenses}

%---------------------------------------------------------------------
\begin{lab}[Interval, Test and Effect Size]
\begin{lstlisting}[language=Python]
import numpy as np
from scipy import stats

old = np.array([0]*32 + [1]*68)      # 68 of 100 passed
new = np.array([0]*21 + [1]*79)      # 79 of 100 passed

# --- confidence interval for the new format's pass rate ---------------
p, n = new.mean(), len(new)
se = np.sqrt(p*(1-p)/n)
print(f"new pass rate: {p:.2f}  95% CI [{p-1.96*se:.3f}, {p+1.96*se:.3f}]")

# --- hypothesis test --------------------------------------------------
t, pval = stats.ttest_ind(new, old)
print(f"t = {t:.2f}, p = {pval:.4f}")

# --- effect size: report this WITH the p-value, never instead of it ---
pooled = np.sqrt((old.std(ddof=1)**2 + new.std(ddof=1)**2) / 2)
print(f"Cohen's d = {(new.mean()-old.mean())/pooled:.3f}")

# --- the multiple-comparisons demonstration ---------------------------
rng = np.random.default_rng(0)
hits = sum(stats.ttest_ind(rng.normal(size=50), rng.normal(size=50)).pvalue < 0.05
           for _ in range(20))
print(f"'significant' results from 20 tests on pure noise: {hits}")
\end{lstlisting}

\textbf{Run the last block a few times.} It is drawing from identical
distributions every time, so every ``significant'' result is false. Seeing this
happen is more convincing than being told it happens.
\end{lab}

%---------------------------------------------------------------------
\begin{checkpoint}
\begin{tightnum}
  \item A 95\% CI for mean response time is $[210, 260]$~ms. Is the statement
        ``there is a 95\% chance the true mean is between 210 and 260'' correct?
        If not, state the correct interpretation.
  \item A study reports $p = 0.001$ with $n = 500{,}000$ and a difference of
        0.03\%. Should the organisation act on it?
  \item You test 15 features for association with an outcome at $\alpha = 0.05$.
        Roughly what is the chance of at least one false positive if none are
        genuinely associated?
\end{tightnum}
\end{checkpoint}

%---------------------------------------------------------------------
\begin{chaptersummary}
\begin{tight}
  \item Inference generalises from sample to population; a biased sample cannot be
        rescued by size.
  \item A confidence interval describes the reliability of the \emph{procedure},
        not the probability of this particular interval.
  \item Precision improves only as $\sqrt{n}$ --- quadruple the data to halve the
        error.
  \item A $p$-value is $P(\text{data this extreme} \mid H_0)$. It is not the
        probability that $H_0$ is true, and it says nothing about effect size.
  \item Type I errors are false alarms, Type II are misses; which is worse is a
        domain judgement.
  \item Report effect size with every $p$-value, and correct for multiple
        comparisons or expect to find things that are not there.
\end{tight}
\end{chaptersummary}

\begin{reviewq}
  \item \textbf{(MCQ)} The $p$-value is the probability of: (a)~$H_0$ being true
        (b)~data at least this extreme given $H_0$ (c)~a Type II error
        (d)~the effect being large.
  \item \textbf{(MCQ)} To halve a confidence interval's width you must multiply the
        sample size by: (a)~2 (b)~4 (c)~$\sqrt{2}$ (d)~10.
  \item \textbf{(Short)} Define Type I and Type II errors and give a security
        example of each with its practical cost.
  \item \textbf{(Short)} Explain why voluntary-response sampling is especially
        dangerous in a study of student disengagement.
  \item \textbf{(Short)} Why does an IoT dataset of ten million readings make
        $p$-values nearly useless?
  \item \textbf{(Applied)} A trial compares two intrusion detection configurations
        over 30 days. Config A raised 220 alerts of which 14 were real; config B
        raised 95 of which 11 were real. Set up the hypotheses, say which test you
        would run, and explain what effect size you would report to management.
  \item \textbf{(Applied)} A colleague tests 40 candidate features against an
        outcome and reports the three with $p < 0.05$. Explain what is wrong,
        compute the Bonferroni-corrected threshold, and describe a defensible
        procedure they should have followed instead.
\end{reviewq}
